\section{More Discussions}
\subsection{(Potential) More Essential Factors}
\noindent\textbf{Accurate understanding of the videos}\quad Through our analysis of failure cases, we've observed that many models struggle with accurately describing videos. While they might be adept at detecting objects within the videos, they often falter in comprehending the contextual relationship between sequential events. Such misinterpretations indicate that there's a need for further exploration in this domain. The videos we've used can indeed serve as an invaluable dataset for probing video descriptions in depth.\\
\noindent\textbf{Logic Reasoning}\quad The primary nature of our videos encompasses content that is counterintuitive and contradicts common sense. For models to understand these, it's imperative they grasp the concept of "common sense." They must deduce what would typically transpire under normal circumstances and then use that perspective to humorously interpret the video. This necessitates the model to possess strong reasoning capabilities, especially when it comes to common sense reasoning.\\
\noindent\textbf{Extra Knowledge - Sense of Humor}\quad To decipher the humor in a video, it's plausible that understanding the fundamental principles of humor is crucial. This type of knowledge, along with many other tidbits of common sense and additional information, might enhance the model's performance. Determining how to integrate valuable knowledge and discerning what counts as "valuable" are topics that warrant further exploration.

\subsection{Potential Solutions}
\noindent\textbf{Model Size}\quad Increasing the number of parameters is a natural method to enhance the model's performance. However, this approach comes with its own set of engineering challenges, requiring improvements in model optimization and deployment. We're also curious about the relationship between the number of parameters and the performance on FunQA. This is an intriguing research point in itself, and our dataset can serve as an excellent test bed to further this exploration.\\
\noindent\textbf{Data Quality}\quad We believe the emphasis for this task should be on data collection. Current trends with large and dynamic models have shown that having vast amounts of low-quality data isn't as effective as a smaller quantity of high-quality data. Thus, we hope the community can discover the type of data that genuinely assists in understanding counterintuitive videos. This is a crucial research direction.\\
\noindent\textbf{Training Strategy}\quad Studying training strategies is also essential. For instance, determining which type of data to start learning from, and understanding the significance of curriculum learning, among others.\\
\noindent\textbf{Model Collaboration}\quad Ultimately, we might not need to solely focus on a single model to solve this problem. Perhaps multiple models collaboratively working on examples in an elegant manner could be a method to enhance performance. However, this approach might necessitate paying more attention to the overall efficiency of model implementation

\subsection{The Emphasis of Temporal Dynamics in FunQA}
\textbf{HumorQA Example}: Humor\_Example.mp4\\
\textbf{H2}: An individual slipped on the staircase filled with ice and tumbled down to the very bottom, followed by a second person who also fell after witnessing the first person's fall.\\
\textbf{H3}: The first person falling down the stairs step by step was already very funny, and the second person repeating the same mistake and falling down made it even more hilarious.\\
\textbf{Explanation}: As answered in H2, the main element of humor in this video is two people slipping down an icy staircase in a sequential order, and if the model does not make sense of the temporal information, it will not be able to give the sequential logic of the two people slipping down one after the other. It is also clear from the answer to H3 that the sequence of slipping backwards and forwards is one of the sources of humor, as the person at the back does not learn from the lesson of the former and slips down in a similar way, and this repetition of the wrongdoing brings humor to the situation.\\
\textbf{CreativeQA Example}: Creative\_Exapmle.mp4\\
\textbf{C2}: At the center of the stage is a blue rectangular box. Following the little girl's watering of the box, four individuals lying inside gradually lift their hands and legs, each at different paces and heights, until they stop and reveal green painted leaf-shaped objects on their limbs.\\
\textbf{C3}: The main creative element of this video is when the four individuals in a box, with varying heights and movement speeds, gradually raise their hands and feet to mimic the growing process of the carrot seeds planted by the little girl. The green leaf-shaped objects, tied and opened in advance with their hands and feet, are used to simulate the true sprouting of carrot sprouts, resulting in a lively and imaginative scene.\\
\textbf{Explanation}: The performance in this video mimics growing bean sprouts, the process of growing beans into sprouts by constantly watering the soil. Trying to describe the content of this video idea and explain where the idea came from requires an understanding of the temporal information. Analyzing the visual information together, the model has to understand the sequence of watering, bean sprouts growing from the soil, and bean sprouts growing taller and greener in order to answer the question accurately. This difference in understanding can be demonstrated on the C3 task, where without analyzing the temporal information, the optimal answer would be bean sprouts growing from the soil, and cause and effect logical relationships such as watering to grow bean sprouts would be ignored.